\section{Eliminating functional unit observations}\label{sec:units}

A unit observation substitutes the number one in every coordinate of a free
functional variable.  It is an observation of the assignment, not a claim
that a closed term has the all-one vector.  We first show how to eliminate
such observations without requiring that inhabitant.

\begin{lemma}[One-name scalar elimination]\label{lem:one-unit}
Let $A:\tau$ be an actual source, $h=\can A$, $m=\lv(\tau)$, and let
$X:\sigma$ be one of its free functional names.  Put
\[
 n=\lv(\sigma)+1,\quad d=2^n\sum_{i=0}^m\sz^X(h_i),\quad
 P=(\partial^X_0h_0)_0,\quad R=h_0\{X:=\mathbf1\},\quad L=\lambda X.A.
\]
The graph of the desired scalar $R$ has a finite positive-coefficient linear
definition using scalar observations of actual sources without free $X$.
\end{lemma}
\begin{proof}
The original abstraction clause gives $L_0=dP+R$.  If the whole $h_0$ is
$X$-free, the stopping rule gives $P=R+1$, hence
$L_0=(d+1)R+d$.  The positive coefficient uniquely determines $R$.

Suppose $h_0$ is active.  Introduce a fresh unused ground $u$ and put
$V=\lambda u:o.A$, $k=2(m+1)$, $p=k+1$, $h'=\can V$,
$d'=2^n\sum_j\sz^X(h'_j)$ and $L'=\lambda X.V$.
The literal scalar is $h'_0=k(h_0+1)+h_0$, where the multiplication means
the original finite sum of copies, before any collecting.  Its partial is
$pP+4k$: a copy of $h_0+1$ contributes $P+3$, and the $k$ sum nodes
joining the copies and retained $h_0$ contribute $k$ more.  The retained
unit value is $pR+k$.  The scalar size is
$p\sz^X(h_0)+3k>\sz^X(h_0)$.

No positive-coordinate size decreases.  For $m\ge1$ the new first field
has two copies of $h_1$ and finite constants, and higher fields have two
copies of their old field.  An entire $X$-free field has size one on both
sides; an active field has at least its original size.  For $m=0$ the new
positive field is one.  Therefore $d'>d$.  With $c=4kd'+k$ we obtain
\[
 L'_0=pd'P+pR+c,\qquad
 pd'L_0+dc=dL'_0+p(d'-d)R.
\]
The coefficient of $R$ is positive.  These two sources have no free $X$,
and every other functional name can still receive its common unit observation.
The intended $R$ satisfies the equation and is its unique solution once the
other observations are fixed.  All coefficients are finite syntactic integers.
\end{proof}

\begin{theorem}[Finite unit-observation elimination]\label{thm:unit-elimination}
For an actual source $A$ and a finite set of all its free functional names,
the absolute scalar after their common unit observation is existentially
definable in $\Esrc$, retaining all original shared ground inputs.
\end{theorem}
\begin{proof}
Apply Lemma~\ref{lem:one-unit} to one name, then recursively to the actual
sources $L,L'$ on the remaining set.  Each branch removes a name; the new
$u$ is ground.  Induction on the number of functional names terminates.
Leaves have only ground free names.  Lemma~\ref{lem:absolute} supplies their
absolute scalar graphs, including leaves of functional output type.

For each fixed ground valuation the intended values supply positive witnesses.
Reverse induction through the positive-coefficient equations proves uniqueness
of every internal scalar.  Encoding raw values as naturals plus one merely
changes fixed additive constants.  The Boolean compiler consumes the resulting
finite system.  No repeated original name has acquired an independent value.
\end{proof}

\section{A complete source grammar with ground binders}\label{sec:ground-binders}

Fix a tracked $F:\sigma_k=o^k\to o$, $k\ge1$.  Contexts of level at most
one are generated by $F$, ground variables, supplied actual $F$-free blocks
of result level at most one, ground-argument application, and abstraction
$\lambda z:o.U$.  Nesting and functional-result ground lambdas are unrestricted.
An opaque block may have arbitrary internal types, lambdas and recursors,
and may capture surrounding ground names.  Other free functional names
receive a single common unit observation.

Write $q(U)=\can U_0$ and $h(U)=\can U_1$ when $U$ is functional.
A scalar collapse in an active scalar branch is annotated by its actual
operator source.  This annotation is retained under scalar copying, ground
unit substitution and ground partials.  It is not an assertion that every
scalar subexpression is itself a source.

\begin{lemma}[Affine heads and their finite surplus]\label{lem:affine-head}
Every functional context has an additive first-field tree with leaves $F_1$
and whole $F$-free positive expressions.  At a selected collapse with head
$h$, its tail is $F$-active whenever $h$ is active.  At common units there is
a finite $\delta(h)\ge0$ such that
\[
 (\partial^F_1h)_1=h^*+\delta(h).
\]
The surplus is zero only at the bare leaf $F_1$; a nonfinite $h^*$ has
$\delta(h)\ge1$.  The output-two head partial is finite.
\end{lemma}
\begin{proof}
Ground application retains the operator's positive fields.  A ground lambda
with ground body has first field $(\partial^z_0q)_1$; the whole-free and
variable cases are one, a sum adds its two outputs and one, and an active
collapse contributes its ground-free positive head plus one and the tail
partial.  For functional body the resulting first field is
$(\partial^z_0q)_1+2h+1$.  These clauses prove additive closure by source
induction, and show that occurrence of $F$ in an operator head implies
occurrence in its scalar.  In a collapse its scalar tail includes that
operator scalar.  Ground partials and ground units preserve this support.

On the additive head tree a whole-free stop has surplus one, a tracked
$F_1$ leaf has surplus zero, and an active sum adds its child surpluses
and one.  Only the bare tracked leaf has surplus zero, and its unit value
is finite.  Output two is one at either leaf and adds child values and one
at an active sum, proving its finiteness.
\end{proof}

\begin{lemma}[Actual unary carriers]\label{lem:carriers}
Each selected nonfinite head has an actual unary carrier with that exact
unit first field.  Its ground captures can be closed to zero without changing
the field, and its actual scalar is retained.
\end{lemma}
\begin{proof}
If its annotated operator is $U:o^r\to o$, replace the free tracked name
throughout $U$ by one fresh same-type $J:\sigma_k$ and apply $r-1$ literal
zero arguments.  This yields $V_U:o\to o$.  Each ground application retains
the first field.  The same common unit observation on $J$ reproduces the
original head; this is global same-type renaming, not a change of the result
type of one head occurrence.  Open captured ground names using their fresh
identities and close them to zero.  The generated-source invariant proves
that their values cannot affect the positive field.  No scalar is reset to one.
\end{proof}

\begin{lemma}[Multiset accounting and ground transport]\label{lem:reserve}
For every annotated scalar $q$ there are an actual-carrier multiset
$\mathcal M(q)$ and a finite $c(q)\ge0$ with
\[
 (\partial^F_0q)_1=\sum_{V\in\mathcal M(q)}V_1+c(q),\qquad
 d(q)=c(q)-3|\mathcal M(q)|.
\]
On the generated active sources, $d\ge1$; it is at least four except
possibly at the bare functional source $F$.  Ground units preserve the
multiset and $d$, and a ground scalar partial preserves the multiset and
does not decrease $d$.
\end{lemma}
\begin{proof}
Define the accounting literally.  At an $F$-free scalar or $F_0$, use the
empty multiset and $d=1$.  At an active sum use multiset union and
$d=d_1+d_2+1$.  At an active collapse $\psi(\omega\otimes h+t)$, the
tail is active by Lemma~\ref{lem:affine-head}.  For finite $h^*$ retain its
tail multiset and add $h^*+\delta(h)$ to $d$.  For nonfinite $h^*$ append
$V_U$ and add $\delta(h)-3$.  The canonical positive partial proves each
identity; subtraction here is only in the finite accounting integer.

A ground substitution changes no positive head and retains active nodes,
so its accounting is identical.  For a ground partial $D_zq$, an $F$-free
piece remains a free stop.  On an active $z$-free $q$ the new $q+1$ adds
two to $d$.  An active sum is transformed into
$(D_zq_1+D_zq_2)+1$, again adding two beyond the child contributions.
At an active collapse the ground-free head and its carrier decision do
not change, and induction handles the tail.  This proves the transport
statements for every nested ground binder.

It remains to check the lower reserve on sources.  The base $F$ has one.
An active ground application has scalar $\psi(\omega\otimes h+q_U+q_V)$.
If $h^*$ is finite, its reserve is
$d(q_U)+d(q_V)+1+h^*+\delta(h)\ge d(q_U)+d(q_V)+2$,
since $h^*\ge1$.  If it is nonfinite, the reserve is
$d(q_U)+d(q_V)+\delta(h)-2\ge d(q_U)+d(q_V)-1$.
An active ground argument is not the bare functional $F$, so has reserve
at least four by induction.  If the argument is inactive, the active operator
cannot be bare $F$ in this nonfinite case, whose first unit field is one.
The reserve is again at least four.

For an active ground lambda with body $U$, the exact scalar is
$vD_zq(U)+q(U)\{z:=1\}$, with the full-vector syntactic multiplier
$v=2\sum_i\sz^z(\can U_i)\ge2$.  There are $v+1$ active copies and
$v$ joining sum nodes.  The transport result gives reserve at least
$(v+1)d(q(U))+v\ge5$.  It includes the positive-field size of a functional
body.  This completes the source induction and proves the claimed bounds.
\end{proof}

\begin{lemma}[One final abstraction]\label{lem:chain}
For actual unary sources $V_j$ free of a fresh $x$,
\[
 \can{\lambda x.V_1(\cdots(V_mx)\cdots)}_1
       =\sum_{j=1}^m(V_{j,1}+3)+1.
\]
\end{lemma}
\begin{proof}
The positive ground partial starts at one on $x$.  One application adds
the operator first field plus one from its whole-free positive partial;
its scalar tail adds two more through its operator scalar and sum node.
Thus it adds $V_{j,1}+3$.  Induction on the list proves the formula, with
one final lambda and no charge for reabstracting each copy.
\end{proof}

\begin{theorem}[Original-budget ground-binder construction]\label{thm:ground-construction}
Let $P:o$ be $F$-active in the complete grammar, $B=\lambda F.P$, and
let $C:\sigma_k$ be an actual $F$-free caller.  Write $b$ for the second
field of the original full box $\can B\square\can C$.
Then $b=B_2$ is finite, $b\ge4$, and at $\mathcal N=2^b$ one effectively
constructs an actual unit-open unary $K$ with
\[
 \mathcal N(B_1+C_1)=K_1+\Delta,\qquad \Delta\in\N,\quad\Delta>0.
\]
No child construction certificate or enlargement of the original scale is
assumed.
\end{theorem}
\begin{proof}
Since $P$ is ground, $B_1$ is its positive scalar partial.  Lemma~\ref{lem:reserve}
gives a multiset and $d\ge4$ such that
$B_1=\sum_{V\in\mathcal M}(V_1+3)+d$.
The output-two partial is finite by Lemma~\ref{lem:affine-head}.  At an
active application that row is the head row plus the operator and argument
scalar-partial rows and one, hence at least four.  At an active lambda the
same copying computation gives at least $2v+1\ge5$.  The sole exception,
bare $F$, is not ground.  Thus $B_2\ge4$; the caller has level one, so
the original untruncated second box is exactly $B_2$.

Unary-saturate the actual caller to $U_C=C0^{k-1}$, retaining $U_{C,1}=C_1$.
In one flat list take $\mathcal N$ copies of every multiset occurrence and
$\mathcal N$ copies of $U_C$.  Lemma~\ref{lem:chain} gives
\[
 K_1=\mathcal N\sum_{V\in\mathcal M}(V_1+3)
             +\mathcal N(C_1+3)+1,
 \quad\Delta=\mathcal N(d-3)-1\ge\mathcal N-1>0.
\]
Every carrier is actual by Lemma~\ref{lem:carriers}.  Replication is finite
and effective, though it may be enormous.  There is no unsupported derivative
realizer and no change to the original caller or its scalar value.
\end{proof}

\section{From a certificate to the whole original head graph}\label{sec:head-graph}

The next result is a consumer with an explicit premise.  It applies to
the preceding theorem, but also to any other source construction paying
the same premise.

\begin{theorem}[Whole-head graph from a paid certificate]\label{thm:head-graph}
Suppose actual sources $B,C,K$ and a finite original box $b$ supply
$a=2^b(B_1+C_1)=K_1+\Delta$ at common functional units, where
$\Delta\ge0$ is finite and $K:o\to o$ is unit open.  Then the complete
genuine graph of
\[
 J_\xi(t)=\psi(\xi+t-\norm(\xi)),\quad
 \xi=\omega\otimes a,\quad t\ge\norm(\xi),
\]
is effectively existentially definable in actual-source arithmetic.  It
retains the original scalar/input wires of $B,C$ and every finite-prefix input.
\end{theorem}
\begin{proof}
Close only auxiliary ground names of $K$ to literal zero, obtaining $K^0$.
Its positive field is unchanged by the source invariant; no ground name in
the original $BC$ is closed.  After functional units set
\[
 k=K^0_1,\quad c=K^0_0,\quad\beta=\omega\otimes k,\quad
 C_\beta=\norm(\beta)+c+1.
\]
These are fixed computable values.  Theorem~\ref{thm:unit-elimination}
observes the exact ray $\can{K^0z}_0=J_\beta(C_\beta+n)$ for $z_0=n+1$,
without replacing its real scalar $c$ or inventing a unit inhabitant.

The certificate gives $\xi=\beta+\omega\Delta$ and
$\norm(\beta+\omega j)=\norm(\beta)+2j$.  At a genuine total norm $t$,
the exact predecessor of $\beta+\omega j$ is
$\beta+\omega(j-1)+\Phi(t)-\norm(\beta+\omega(j-1))$.
The maximal-predecessor theorem of \cite{wwfoundations} proves this by
retaining the affordable prefix and spending the rest on the finite tail.
Consequently induction on $j$ gives the \emph{upward} identity
\[
 J_\xi(t)=J_\beta(\Phi^\Delta(t))+\Delta
          \qquad(t\ge\norm(\xi)).
\]
This does not invert the sparse function $\Phi^\Delta$.

Set $T=\max(\norm(\xi),C_\beta)$.  For $t\ge T$ use the fixed finite
chain of graph equations $u_0=t$, $u_{j+1}=\Phi(u_j)$, and
$u_\Delta=\norm(\beta)+c+z_0$.  Its required $z_0$ is positive because
$u_\Delta\ge t\ge C_\beta$.  Observe $v=\can{K^0z}_0$ through the
unit-elimination theorem and put $y=v+\Delta$.  The identity proves both
existence and uniqueness of $y=J_\xi(t)$ on this tail.

Every $\Phi$ gate is definable by the actual $S$ graph:
$F_S(n)=\Phi(n+4)$, with the four smaller inputs filled by computed
constants.  There is a fixed number $\Delta$ of gates, not an assumed
variable-iteration primitive.  For each genuine $t_0<T$, use the exact
closed evaluator to compute $j_0=J_\xi(t_0)$ and include the disjunct
$t=t_0\wedge y=j_0$.  These finitely many points and the tail define
the entire genuine graph.  Range constraints for the ray stay inside
the tail branch; imposing them globally would lose small inputs.

The Boolean compiler handles this finite disjunction.  Total observation
graphs may be flattened outside it, but branch-specific range constraints
may not.  For the actual original target the shared raw wire is
$t=\norm(\xi)+B_0+C_0+\lv(\sigma)$; the unit-elimination theorem observes
these scalars without closing their original ground names.  Inserting this
numerical wire changes no original source or valuation.  All closed
evaluations, replication, literal lengths and finite-prefix enumeration
are charged to effectivity, with no feasible complexity assertion.
\end{proof}

The grammar restriction is real.  With $\tau=o\to o$ and
$H=R^o0(D^o0):\tau$, one has $H_1=2\omega$.  For $F:\tau\to\tau$,
the first field of $FH$ is $h=2^{F_2}\otimes(F_1+2\omega)$.
Its unit value is $4\omega+2$, while its positive $F$ partial is
$2\omega+6$.  No finite reserve makes these equal.  Thus an active
functional argument can defeat the affine invariant; the example does
not refute every possible source-construction method.
